\chapter{Laumon 因子化}\label{ch:7}\index{Laumon 因子化|hyperpage}
% 原PDF第40页中部
\section{洁净性中剩下的引理}\label{sec:7.1}
本章完成“消失推出洁净性”的证明的最后一步。固定 $M_k\in\Bun_k^0$、$J\in\Tor^\ell$、$\mathring M\in\Bun_k$ 以及 $a:M_k\to J$，其中 $\Bun_k^0$ 表示作为 Whittaker 旗终点出现的归一化次数零分支。假设 $\mathring M\oplus J\in C_k$。考虑叠
\[H:=\left\{\begin{array}{l}0\to\mathring M\to M\to J\to0,\quad M\in C_k,\\u:M_k\hookrightarrow M,\quad(M_k\xrightarrow{u}M\twoheadrightarrow J)=a\end{array}\right\}.\]
它有映射 $q_H:H\to\Tor$，$(M,u)\mapsto M/M_k$。洁净性证明所需的引理是紧支撑消失
\[\RGamma_c(H,q_H^*\cL_E)\simeq0.\]

还要对良好的开子叠 $C_k$ 作一个技术性加强。可以要求：对每个扩张 $0\to\mathring M\to M\to J\to0$，
\[\mathring M\oplus J\in C_k\quad\Longrightarrow\quad M\in C_k.\]
确保这一蕴涵的一种方法，是选取线丛 $L^{\mathrm{est}}$ 并施加额外的开条件
\[\Hom(M,L^{\mathrm{est}})=0.\]
事实上，对挠层 $J$ 和无挠目标 $L^{\mathrm{est}}$，$\Hom(J,L^{\mathrm{est}})=0$。因此 $\mathring M\oplus J\in C_k$ 等价于 $\mathring M\in C_k$，对扩张应用 $\Hom(-,L^{\mathrm{est}})$ 得到 $M\in C_k$。于是，在纤维计算中可以忽略条件 $M\in C_k$。
\section{满射情形}\label{sec:7.2}
\begin{proof}
首先假设 $a:M_k\twoheadrightarrow J$。置 $\mathring M_k:=\ker(a)$。则 $H$ 的一个对象等同于如下图：
\origdiagram{p040-exact.pdf}{0.4645}
等价地，它是一个修改 $\mathring M_k\hookrightarrow\mathring M$。因此
\[H\simeq\{\mathring M_k\}\times_{\Bun_k}\operatorname{Mod}_k^d\times_{\Bun_k}\{\mathring M\},\qquad\Bun_k\longleftarrow\operatorname{Mod}_k^d\longrightarrow\Bun_k.\]
在这一等同下，$q_H$ 记录 $\mathring M/\mathring M_k$。忽略 $\Av_E$ 定义中无碍于论证的归一化一维线后，
\[\RGamma_c(H,q_H^*\cL_E)\simeq(\Av_E^d(\delta_{\mathring M_k}))_{\mathring M}\otimes(\text{一维线}),\qquad d=\ndeg\mathring M-\ndeg\mathring M_k.\]
\srcpage{41}
在洁净性证明所用的归一化分支中，$\ndeg M_k=0$，所以 $\ndeg\mathring M_k=-\ell$。条件 $\mathring M\oplus J\in C_k$ 给出 $\ndeg\mathring M>nk(2g-2)$。故
\[d=\ndeg\mathring M-\ndeg\mathring M_k>nk(2g-2)+\ell>nk(2g-2).\]
因此，$\Av_E^d$ 的消失定理给出 $\RGamma_c(H,q_H^*\cL_E)=0$。
\end{proof}
\section{归约到 \texorpdfstring{$M_k\to J$}{Mₖ→J} 的像}\label{sec:7.3}
\begin{proof}
对一般映射 $a:M_k\to J$，令 $J'\subset J$ 为其像，置 $J'':=J/J'$。令 $\mathring M_k=\ker(M_k\twoheadrightarrow J')$。对 $H$ 的一个对象，置 $M':=M\times_JJ'$。于是有
\origdiagram{p041-three.pdf}{0.4690}
令 $H'$ 为前两行构成的叠；它是以 $J'$ 为目标的满射版本。

令 $\Tor^{\mathrm{ses}}$ 为挠层短正合列
\[0\longrightarrow J_s\longrightarrow J_m\longrightarrow J_q\longrightarrow0\]
的叠。记
\[m(J_\bullet)=J_m,\qquad s(J_\bullet)=J_s,\qquad q(J_\bullet)=J_q.\]
商序列
\[0\longrightarrow M'/M_k\longrightarrow M/M_k\longrightarrow M/M'\longrightarrow0\]
定义映射 $H\to\Tor^{\mathrm{ses}}$。这里
\[M'/M_k\simeq\mathring M/\mathring M_k,\qquad M/M'\simeq J/J'.\]
诱导图
\origdiagram{p041-diagram.pdf}{0.2914}
不是笛卡尔图。然而，对于所考虑的层，基变换仍然成立。原因在于
\[H\longrightarrow\Tor^{\mathrm{ses}}\times_{\Tor\times\Tor}H'\]
是仿射丛。在由
\[0\to J_s\to J_m\to J_q\to0,\qquad0\to M_k\to M'\to J_s\to0\]
分类的点处，纤维是仿射空间
\[\Ext^1(J_m,M_k)\times_{\Ext^1(J_s,M_k)}\{[0\to M_k\to M'\to J_s\to0]\}.\]
\srcpage{42}
其维数在每个连通分支上恒定；事实上，这一仿射空间以
\[\ker(\Ext^1(J_m,M_k)\to\Ext^1(J_s,M_k))\simeq\Ext^1(J_q,M_k)\]
为平移空间，维数为 $k\length(J_q)$。所以，忽略恒定的仿射丛移位与 Tate 扭转后，只需证明
\[\RGamma_c((s,q)_!m^*\cL_E|_{H'})=0.\]
由满射情形，这将由下面的因子化恒等式推出。
\end{proof}
\section{Laumon 层通过短正合列的因子化}\label{sec:7.4}
\begin{proposition}{Laumon 因子化\cite{KL85}}{7.4.1}
对任何能够定义 Laumon 复形的局部系统 $E$，特别是洁净性论证所用的几何不可约 $E$，在 $\Tor\times\Tor$ 上有
\[(s,q)_!m^*\cL_E\simeq\cL_E\boxtimes\cL_E.\]
\begin{transnote}{含糊点拨}
这条恒等式就是「因子化」的层化陈述：$\cL_E$ 是沿 Hecke 对应把 $E$ 拉回、再把修改方向的层推前所得的核，而 $(s,q)_!m^*\cL_E$ 把 $\Tor\times\Tor$ 上的乘法拆开，恰好对应「把挠层分成两个子挠层、再各自配 $E$」这个操作。命题 7.4.1 把它当作输入，证明在 \chref{9} 的 Laumon 层构造与 §7.1 的 Laumon 复形里完成；「能够定义 Laumon 复形」只要求所选层理论合适，几何不可约是洁净性论证另加的条件，两者不可混为一谈。[依据：原书 p.42 命题 7.4.1 及其前的构造；并见 [KL85]。]
\end{transnote}
等价地，在 $\Tor^{\mathrm{ses}}$ 中子对象长度为 $d_1$、商长度为 $d_2$ 的分支上，在 $\Tor^{d_1}\times\Tor^{d_2}$ 上有
\[(s,q)_!m^*\cL_E^{d_1+d_2}\simeq\cL_E^{d_1}\boxtimes\cL_E^{d_2},\]
其中移位与 Tate 扭转已经包含在 $\cL_E$ 的归一化中。
\end{proposition}
\begin{note}\numberednote{7.4.2}
前后的洁净性论证假设 $E$ 几何不可约，但因子化恒等式本身是 Springer–Laumon 构造的结构性质，不使用这种不可约性。
\end{note}
\begin{proof}
回顾正则半单轨迹
\[\Tor^{d,\mathrm{rss}}=\left\{\bigoplus_{i=1}^dk(x_i)\ \middle|\ x_i\ne x_j\text{ 当 }i\ne j\right\}\subset\Tor^d,\]
它是 $X^{(d),\mathrm{rss}}=X^{(d)}-\Delta$ 上的胚叠，其在几何点处的稳定子为 $(\Gm)^d$。
\begin{transnote}{含糊点拨}
胚叠是指局部同构于 $BG$ 的叠；这里每个 $k(x_i)$ 上的标量自同构各给出一个 $\Gm$，所以稳定子是 $(\Gm)^d$，而商叠局部是 $B(\Gm^d)$。这一步是为下一行的对称下降铺路：正因为 $\Tor^{d,\mathrm{rss}}$ 相对于 $X^{(d),\mathrm{rss}}$ 只剩 $(\Gm)^d$ 的自同构，$E^{\boxtimes d}$ 沿有序支撑覆盖的推前才带 $S_d$-作用，于是可以取 $S_d$-不变部分把它下降到 $X^{(d),\mathrm{rss}}$ 上。[来源：nLab "gerbe"，https://ncatlab.org/nlab/show/gerbe（Last revised 2026-02-14；2026-11-17 访问）；并见原书 p.42 本段与第 4 章对 $(\Gm)^d$-胚叠的讨论。]
\end{transnote}若 $r_{\mathrm{rss}}:X^{d,\mathrm{rss}}\to X^{(d),\mathrm{rss}}$ 为有序支撑覆盖，$s_{\mathrm{rss}}:\Tor^{d,\mathrm{rss}}\to X^{(d),\mathrm{rss}}$ 为支撑映射，则
\[E^{(d)}:=(r_{\mathrm{rss},*}E^{\boxtimes d})^{S_d},\qquad\cL_E^d\simeq j_{\mathrm{rss},!*}(s_{\mathrm{rss}}^*E^{(d)}[d]).\]
\begin{transnote}{含糊点拨}
$j_{\mathrm{rss},!*}$ 是中间延拓（IC 延拓）：对开嵌入 $j_{\mathrm{rss}}$，它是像 $j_!\to j_*$ 所对应的那个中间层，既非紧支撑延拓 $j_!$，也非 $j_*$。取它才得到 Goresky–MacPherson 延拓，与 §9.1 定义的 Laumon 层一致。$[d]$ 把 $E^{(d)}$ 的度数移到 $\Tor^d$ 的维数上，使 $\cL_E^d$ 落在反常 $t$-结构的度 0 层；若忘记这个移位，下一章的洁净性就会差一个平移。[依据：原书 p.42 本式与 §9.1；并见 nLab "intermediate extension" 的中间延拓定义（该条目无 Last revised，2026-11-17 访问）。]
\end{transnote}
令
\[U_{d_1,d_2}:=(\Tor^{d_1,\mathrm{rss}}\times\Tor^{d_2,\mathrm{rss}})_{\mathrm{disj}}\]
为两个支撑互不相交的开轨迹，令 $s_i:\Tor^{d_i,\mathrm{rss}}\to X^{(d_i),\mathrm{rss}}$ 为两个支撑映射。在这一轨迹上，两个挠层之间的每条短正合列都唯一分裂，中间项映射 $m$ 等同于除子的加法。因此
\[(s,q)_!m^*\cL_E^d|_{U_{d_1,d_2}}\simeq\operatorname{add}^*s_{\mathrm{rss}}^*E^{(d)}[d]|_{U_{d_1,d_2}},\]
而
\[(\cL_E^{d_1}\boxtimes\cL_E^{d_2})|_{U_{d_1,d_2}}\simeq(s_1^*E^{(d_1)}[d_1]\boxtimes s_2^*E^{(d_2)}[d_2])|_{U_{d_1,d_2}}.\]
这两个局部系统通过对称下降典范等同。此外，右边是其到 $U_{d_1,d_2}$ 的限制的 IC 延拓：乘积局部系统已经光滑地跨越两个支撑相遇的交叉对角线延拓。因此，只需证明
\[(s,q)_!m^*\cL_E^d\simeq\operatorname{IC}((s,q)_!m^*\cL_E^d|_{U_{d_1,d_2}}).\]

由于 $\cL_E^d=(\operatorname{Spr}_E^d)^{S_d}$ 是 Springer–Laumon 层的直和项，只需对 $\operatorname{Spr}_E^d$ 证明相应陈述，使其对 $S_{d_1}\times S_{d_2}\subset S_d$ 等变，再取
\srcpage{43}
不变直和项。该断言在 $X$ 上是 étale 局部的。作这样的局部化后，局部系统 $E$ 成为纤维为 $V$ 的常值系统；系数因子 $V^{\otimes d}$ 随整个构造带入。因此，几何输入就是常系数的 Springer 陈述。

对 $X=\A^1$，长度 $d$ 挠层是带有一个自同态的 $d$ 维向量空间，因此
\[\Tor^d\simeq\mathfrak{gl}_d/\GL_d.\]
\begin{transnote}{含糊点拨}
$\A^1=\Spec k[x]$ 上长度 $d$ 的挠层就是一个 $d$ 维 $k[x]$-模：$x$ 的作用给出一个 $d\times d$ 矩阵，模同构就是 $\GL_d$-共轭，于是模叠恰是 $\mathfrak{gl}_d/\GL_d$，全空间都允许，不额外要求 $x$ 幂零。$x$ 幂零的那些点构成其中的幂零锥 $\cN_d$，正是第四章局部模型在 $0$ 附近看到的部分；这里取整块 $\mathfrak{gl}_d$，才能让 $\Tor^d$ 在 $0$ 处带出幂零锥。[来源：nLab "coherent sheaf"，https://ncatlab.org/nlab/show/coherent+sheaf（Last revised 2026-01-25；2026-11-17 访问）；并见原书 p.43 本段。]
\end{transnote}
对 $d=d_1+d_2$，叠 $\Tor^{\mathrm{ses},d_1,d_2}$ 是抛物对应
\[\mathfrak g/G\longleftarrow\mathfrak p/P\longrightarrow\mathfrak m/M,\]
其中 $G=\GL_d$，$P=P_{d_1,d_2}=UM$，$M\simeq\GL_{d_1}\times\GL_{d_2}$，$\mathfrak p$、$\mathfrak m$ 为相应 Lie 代数。所需几何输入是 Grothendieck–Springer 层的标准抛物限制定理\cite{AHJR14,GPR20}：复形
\[(\mathfrak p/P\to\mathfrak m/M)_!(\mathfrak p/P\to\mathfrak g/G)^*\operatorname{Spr}_{\mathfrak g}\]
是从 $\mathfrak m/M$ 中两个特征多项式正则且互不相交的轨迹作出的 IC 延拓，其 Weyl 局部系统由将 $S_d$ 单值作用限制到 $S_{d_1}\times S_{d_2}$ 得到。

作为参考，Bruhat 证明将
\[\widetilde{\mathfrak p}:=\{(x,gB)\mid x\in\mathfrak p\cap\operatorname{Ad}_g\mathfrak b\}\]
按双陪集部分 $PwB/B$ 分层，其中 $w\in W_M\backslash W_G$。对最小陪集代表 $w$，到 $\mathfrak m$ 的 Grothendieck–Springer 解消的映射是仿射纤维化，其纤维为
\[\mathfrak u_w:=\mathfrak u\cap\operatorname{Ad}_w\mathfrak b,\]
而非整个 $\mathfrak u$；维数依赖于 $w$。采用通常的反常归一化，这些分层组合成上述抛物限制的 IC 延拓，不会出现支撑在正则且不交轨迹的补集上的直和项。

带着系数 $V^{\otimes d}$ 取 $S_d$ 不变直和项，将这一 IC 延拓等同于 $\cL_E^{d_1}\boxtimes\cL_E^{d_2}$，因为在正则且不交轨迹上，它恰好就是已描述的对称下降等同。这证明了 Laumon 因子化命题。
\end{proof}
\section{完成引理}\label{sec:7.5}
\begin{proof}
将命题\ref{pro:7.4.1}应用于上文所得剩余表达式：
\[(s,q)_!m^*\cL_E|_{H'}\simeq(\cL_E\boxtimes\cL_E)|_{H'}.\]
第一个因子恰好是 $M_k\twoheadrightarrow J'$ 的满射情形，第二个因子是固定的剩余商 $J/J'$。所以，由满射情形，沿 $H'$ 的 $!$ 推前消失。再由上述仿射丛基变换，
\[\RGamma_c(H,q_H^*\cL_E)=0.\]
这证明了剩下的引理，并完成洁净性证明中缺少的一步。
\end{proof}
